% ID: SEL-D0034-OA-E01
% ESTADO: cotejada. Segunda pasada visual y borrador compilado revisados.
% Fuente: Selectividad/matematicas_aplicadas/titular_septiembre_1_MATEMATICAS_APLICADAS.pdf
% SHA-256: c637bbde5c26c0b35b59202bd832ddc97990d3845365ce483fbf2a49832f73fe
% Página: 1; opción A; ejercicio original 1.

Sean las matrices $A=\begin{pmatrix}1&2\\-1&-3\end{pmatrix}$, $B=\begin{pmatrix}2&-1&3\\4&0&1\end{pmatrix}$ y $C=\begin{pmatrix}-1&1&0\\2&3&-2\end{pmatrix}$.
\begin{enumerate}[label=\alph*)]
\item \textbf{(1.7 puntos)} Resuelva la ecuación matricial $A^2\cdot X+C=2B$.
\item \textbf{(0.8 puntos)} \textquestiondown{}Qué dimensiones deben tener las matrices $P$ y $Q$ para que las matrices $(B+C)\cdot P$ y $B\cdot Q\cdot C^t$ sean cuadradas?
\end{enumerate}
